\begin{lemma}[A common in-measure limit for the complex Riesz images of $L^1$ approximating sequences]
Let $f \in L^1(\gamma_n;\mathbb{C})$.  Then there exist an element $g$ of
$\mathcal{M}(\gamma_n;\mathbb{C}^n)$, the space of $\gamma_n$-almost-everywhere
equivalence classes of measurable $\mathbb{C}^n$-valued functions on
$\mathbb{R}^n$, and a sequence $(v_k)_{k\in\mathbb{N}}$ in
$L^2(\gamma_n;\mathbb{C})$, such that
\begin{enumerate}
\item $\|\iota_{\mathbb{C}} v_k\|_{L^1(\gamma_n;\mathbb{C})}
  \leq \|f\|_{L^1(\gamma_n;\mathbb{C})}$ for every $k$, with $\iota_{\mathbb{C}}$
  the inclusion of \noderef{ComplexToL1};
\item $\iota_{\mathbb{C}} v_k \to f$ in $L^1(\gamma_n;\mathbb{C})$;
\item $R_\gamma^{\mathbb{C}} v_k \to g$ pointwise $\gamma_n$-a.e., with
  $R_\gamma^{\mathbb{C}}$ as in \noderef{ComplexRieszTransform}; and
\item for \emph{every} sequence $(w_k)$ in $L^2(\gamma_n;\mathbb{C})$ with
  $\iota_{\mathbb{C}} w_k \to f$ in $L^1(\gamma_n;\mathbb{C})$, the sequence
  $R_\gamma^{\mathbb{C}} w_k$ converges to $g$ in measure, i.e.
  $\gamma_n\bigl(\{x : \varepsilon \leq |R_\gamma^{\mathbb{C}} w_k(x)
    - g(x)|_{\mathbb{C}^n}\}\bigr) \to 0$ as $k \to \infty$, for every
  $\varepsilon > 0$.
\end{enumerate}

This is the complex counterpart of the Cauchy-in-measure and a.e.-limit
construction that the reference paper invokes in the last sentence of its proof
of the complex-valued estimate, where it says that approximation by complex
$L^2$ functions, with the Cauchy-in-measure and Fatou arguments of its $L^1$
extension subsection, gives the extension to $L^1$.  Clause (4) is the part that
does the real work downstream: it says the limit $g$ depends only on $f$ and not
on the approximating sequence, which is what makes an extension operator
well defined and what forces it to be linear.
\end{lemma}
\begin{proof}
Throughout, $\iota_{\mathbb{C}}$ is the inclusion of \noderef{ComplexToL1}, and
we use from that node that $\iota_{\mathbb{C}} v$ and $v$ have the same
representative, so that $\|\iota_{\mathbb{C}} v\|_{L^1} = \int |v| \, d\gamma_n$,
and that $\iota_{\mathbb{C}}$ is $\mathbb{C}$-linear.

\emph{Part A: the Cauchy-in-measure estimate.}
Let $v, w \in L^2(\gamma_n;\mathbb{C})$ and $\varepsilon > 0$.  By clauses (1)
and (2) of \noderef{ComplexRieszTransformLinear}, applied with the scalar
$c = -1$ for the second,
\[
  R_\gamma^{\mathbb{C}}(v - w)
    = R_\gamma^{\mathbb{C}} v + R_\gamma^{\mathbb{C}}\bigl((-1)w\bigr)
    = R_\gamma^{\mathbb{C}} v - R_\gamma^{\mathbb{C}} w .
\]
Apply \noderef{ComplexMainBound} to the complex $L^2$ function $v - w$ at the
level $\varepsilon/2 > 0$:
\[
  \tfrac{\varepsilon}{2}\,
    \gamma_n\bigl(\{x : \tfrac{\varepsilon}{2}
      < |R_\gamma^{\mathbb{C}}(v-w)(x)|_{\mathbb{C}^n}\}\bigr)
    \leq \frac{8}{\pi}\int |(v-w)(x)| \, d\gamma_n(x)
    = \frac{8}{\pi}\,\|\iota_{\mathbb{C}} v - \iota_{\mathbb{C}} w\|_{L^1} ,
\]
the last equality because $\iota_{\mathbb{C}}(v-w)
 = \iota_{\mathbb{C}} v - \iota_{\mathbb{C}} w$ by $\mathbb{C}$-linearity and
because the $L^1$ norm of that class is the integral of $|v-w|$.  Since
$\varepsilon/2 < \varepsilon$ and $R_\gamma^{\mathbb{C}}(v-w)$ agrees
$\gamma_n$-a.e.\ with $R_\gamma^{\mathbb{C}} v - R_\gamma^{\mathbb{C}} w$, the
set $\{x : \varepsilon \leq |R_\gamma^{\mathbb{C}} v(x)
 - R_\gamma^{\mathbb{C}} w(x)|\}$ is contained, up to a $\gamma_n$-null set, in
$\{x : \varepsilon/2 < |R_\gamma^{\mathbb{C}}(v-w)(x)|\}$; monotonicity of
$\gamma_n$ up to null sets (shared context) and division by $\varepsilon/2 > 0$
give
\begin{equation*}
  \gamma_n\bigl(\{x : \varepsilon \leq |R_\gamma^{\mathbb{C}} v(x)
      - R_\gamma^{\mathbb{C}} w(x)|_{\mathbb{C}^n}\}\bigr)
    \leq \frac{16}{\pi}\cdot
      \frac{\|\iota_{\mathbb{C}} v - \iota_{\mathbb{C}} w\|_{L^1}}{\varepsilon} .
  \tag{A}
\end{equation*}

\emph{Part B: the radial truncations.}
Fix a representative $F : \mathbb{R}^n \to \mathbb{C}$ of $f$.  For
$k \in \mathbb{N}$ define $\varphi_k : \mathbb{C} \to \mathbb{C}$ by
\[
  \varphi_k(z) = \begin{cases} z, & |z| \leq k, \\[2pt]
    k\,z/|z|, & |z| > k, \end{cases}
\]
and put $\tau_k = \varphi_k \circ F$.  (This is the complex analogue of the real
two-sided truncation $\max\{-k,\min\{\cdot,k\}\}$: it cuts the modulus at $k$ and
leaves the argument alone.)

\emph{(B1) $\varphi_k$ is continuous, so $\tau_k$ is measurable.}  On
$\{|z| \leq k\}$ and on $\{|z| \geq k\}$ the two formulas are continuous --- the
second because $z \mapsto |z|$ is continuous and nonzero there, as $k \geq 0$ and
$|z| \geq k$ with $|z| > 0$ whenever $k > 0$; for $k = 0$ the second branch is
the constant $0$ --- and on the overlap $\{|z| = k\}$ they agree, since
$k z/|z| = z$ there.  Hence $\varphi_k$ is continuous on $\mathbb{C}$ by gluing
on two closed sets, so Borel measurable, and $\tau_k = \varphi_k \circ F$ is
measurable.

\emph{(B2) the two pointwise bounds.}  For every $z$, $|\varphi_k(z)|
= \min\{|z|, k\}$: if $|z| \leq k$ this is $|z|$, and otherwise it is
$|k z/|z|| = k$.  Hence for every $x$
\begin{equation*}
  |\tau_k(x)| \leq k \qquad\text{and}\qquad |\tau_k(x)| \leq |F(x)| . \tag{B}
\end{equation*}

\emph{(B3) $\tau_k$ lies in $L^2$, and conclusion (1).}  Since $\gamma_n$ is a
probability measure (\noderef{GaussIsProb}), the bound $|\tau_k| \leq k$ makes
$\tau_k$ square $\gamma_n$-integrable; let $\tilde v_k \in L^2(\gamma_n;\mathbb{C})$
be its class.  The second bound in (B) and monotonicity of the integral give
\[
  \|\iota_{\mathbb{C}}\tilde v_k\|_{L^1} = \int|\tau_k|\,d\gamma_n
    \leq \int|F|\,d\gamma_n = \|f\|_{L^1} ,
\]
which is conclusion (1) for this sequence.

\emph{(B4) $L^1$ convergence, i.e.\ conclusion (2).}  For every $x$, as soon as
$k \geq |F(x)|$ we have $\varphi_k(F(x)) = F(x)$, so $\tau_k(x) \to F(x)$, the
sequence being eventually constant at each point.  By (B),
$|\tau_k(x) - F(x)| \leq |\tau_k(x)| + |F(x)| \leq 2|F(x)|$, and $2|F|$ is
$\gamma_n$-integrable because $f \in L^1(\gamma_n;\mathbb{C})$.  Dominated
convergence (shared context) therefore gives
$\|\iota_{\mathbb{C}}\tilde v_k - f\|_{L^1} = \int|\tau_k - F|\,d\gamma_n \to 0$.

\emph{Part C: a rapidly converging subsequence.}
By (B4) we may choose, for each $j \in \mathbb{N}$, an index $\sigma(j)$ with
$\|\iota_{\mathbb{C}}\tilde v_{\sigma(j)} - f\|_{L^1} \leq \tfrac12 8^{-j}$.  Put
$v_j = \tilde v_{\sigma(j)}$.  Conclusion (1) is inherited from (B3), and
$\|\iota_{\mathbb{C}} v_j - f\|_{L^1} \leq \tfrac12 8^{-j} \to 0$ is conclusion
(2) for $(v_j)$.  By the triangle inequality and $8^{-(j+1)} \leq 8^{-j}$,
\begin{equation*}
  \|\iota_{\mathbb{C}} v_j - \iota_{\mathbb{C}} v_{j+1}\|_{L^1}
    \leq \tfrac12 8^{-j} + \tfrac12 8^{-(j+1)} \leq 8^{-j} . \tag{C}
\end{equation*}
Parts B and C together are \noderef{ComplexL1ApproxByL2Truncations}, which states
exactly clauses (1), (2) and the geometric decay $(C)$ for any $f \in L^1(\gamma_n;\mathbb{C})$,
taking as its one hypothesis the norm identity $\|\iota_{\mathbb{C}} u\|_{L^1} = \int|u|\,d\gamma_n$
that Part A supplies; it is what the formalization invokes here.

\emph{Parts D and E: conclusions (3) and (4).}
The sequence $(v_k)$ now satisfies the three hypotheses of
\noderef{ComplexRieszTransformAELimitUnique}: its hypothesis (A) is the
Cauchy-in-measure estimate $(A)$ of Part A, applied to arbitrary $u,w$ and
$\varepsilon > 0$; its hypothesis (L) is conclusion (2), established in Part B; and
its hypothesis (C) is the geometric decay $(C)$ of Part C.  That node therefore
supplies an element $g$ of $\mathcal{M}(\gamma_n;\mathbb{C}^n)$ such that
$R_\gamma^{\mathbb{C}} v_k \to g$ pointwise $\gamma_n$-a.e., which is conclusion
(3), and such that for every sequence $(w_k)$ in $L^2(\gamma_n;\mathbb{C})$ with
$\iota_{\mathbb{C}} w_k \to f$ in $L^1(\gamma_n;\mathbb{C})$ the images
$R_\gamma^{\mathbb{C}} w_k$ converge to $g$ in measure, which is conclusion (4).
Together with (1) and (2) from Parts B and C this proves the lemma.
\end{proof}
